% =========================================================
% CHAPTER 5: SYSTEM IMPLEMENTATION
% =========================================================

\chapter{SYSTEM IMPLEMENTATION}
\label{chap:implementation}

QueryPilot is implemented entirely in Python using a set of open-source
libraries. The database connector is mysql-connector-python (version 8.x),
chosen for direct MySQL 8.0 compatibility without the overhead of an
object-relational mapping layer. Machine learning components use scikit-learn,
a widely adopted Python library for classical machine learning algorithms.
Data processing uses pandas and numpy. The monitoring dashboard uses Streamlit,
a Python-native framework for building interactive web applications from
scripts. Visualization within the dashboard uses matplotlib. The development
environment is Windows 10/11 with Python 3.10 and Visual Studio Code as the
integrated development environment.

The project is organized into six primary packages: database, collector, ml,
sage, workload, and experiments. Each package is implemented as a Python module
with a clear single responsibility. Configuration is managed through environment
variables loaded from a \texttt{.env} file using the python-dotenv library,
ensuring that database credentials are never hardcoded in source files.

\section{MODULE DESCRIPTION}
\label{sec:module_description}

\subsection{Database Connection Module}
\label{subsec:db_connection}

The Database Connection Module, implemented in \texttt{database/connection.py},
provides a single function \texttt{get\_connection()} that reads database
credentials from environment variables and returns a live
\texttt{mysql.connector.CMySQLConnection} object. The module uses the
\texttt{python-dotenv} library to load the \texttt{.env} file at import time,
exposing the host, port, username, password, and database name as environment
variables. This design ensures that all other modules can obtain a fresh
connection with a single import and function call, without duplicating
connection logic. The \texttt{database/load\_tpch.py} script uses this
module to load TPC-H data using MySQL's LOAD DATA LOCAL INFILE command after
enabling the \texttt{local\_infile} global variable.

\subsection{Feature Extraction Module}
\label{subsec:feature_extraction}

The Feature Extraction Module, implemented in
\texttt{collector/feature\_extractor.py}, derives a fixed eight-dimensional
feature vector from any incoming SQL query. The module first performs lexical
analysis on the query string to extract the query type, the number of explicit
JOIN clauses, and boolean flags for WHERE, GROUP BY, and ORDER BY. It then
executes an EXPLAIN query against MySQL and parses the resulting rows to
extract the number of tables in the plan, the sum of estimated rows, and
whether any table uses an index. The EXPLAIN output is parsed by inspecting
the column names from the cursor description, making the extraction robust to
minor variations in MySQL's EXPLAIN output format across versions.

\subsection{Candidate Optimization Generator Module}
\label{subsec:candidate_generator}

The Candidate Generator Module, implemented in
\texttt{collector/candidate\_generator.py}, produces a list of candidate
SQL variants for a given query. The first candidate is always the native
query with no modifications. Subsequent candidates inject MySQL optimizer
hints by inserting a hint comment immediately after the SELECT keyword using
a regular expression substitution. The module first calls EXPLAIN to
identify the first table in the plan, which is used as the hint target. For
join queries, HASH\_JOIN and NO\_BKA hints are also generated using the first
joined table name extracted from the query string. The candidate list is
capped at four entries.

\subsection{Performance Prediction Model Module}
\label{subsec:prediction_model}

The Performance Prediction Model Module, implemented in
\texttt{ml/performance\_model.py}, maintains an ensemble of three
independently trained scikit-learn regressors. The feature order used by the
model is fixed as: query\_type\_encoded, num\_joins, num\_tables, has\_where,
has\_group\_by, has\_order\_by, estimated\_rows, and uses\_index. The
\texttt{features\_to\_vector} function converts a feature dictionary to a
numpy array using this fixed order.

The ensemble training algorithm is as follows:

\begin{enumerate}
  \item Collect (feature vector, execution time) pairs in the training buffer.
  \item When the buffer size reaches the minimum evidence threshold (10), extract
        arrays $X$ (feature matrix) and $y$ (execution times).
  \item Fit each of the three models independently on the same $(X, y)$ data.
  \item Set the \texttt{is\_trained} flag to True.
\end{enumerate}

The \texttt{predict\_with\_uncertainty} method returns the tuple
$(\mu(x), \sigma(x))$ where $\mu(x)$ is the mean of the three predictions
and $\sigma(x)$ is their standard deviation.

\subsection{Uncertainty Estimator Module}
\label{subsec:uncertainty}

The Uncertainty Estimator is implemented inline within the
\texttt{PerformanceModel} class's \texttt{predict\_with\_uncertainty} method.
For a given feature vector $x$, the three model predictions
$\hat{y}_1, \hat{y}_2, \hat{y}_3$ are computed independently. The relative
uncertainty is defined as:

\begin{equation}
  u(x) = \frac{\sigma(x)}{\mu(x)}
  \label{eq:uncertainty}
\end{equation}

where $\mu(x) = \frac{1}{3}\sum_{i=1}^{3}\hat{y}_i$ and
$\sigma(x) = \sqrt{\frac{1}{3}\sum_{i=1}^{3}(\hat{y}_i - \mu(x))^2}$.

A candidate is excluded from SAGE's consideration if its $u(x) > \tau_u = 0.8$.
This threshold was chosen to allow moderate ensemble disagreement while
filtering out candidates where the three models give substantially different
predictions, which typically indicates that the feature vector is unlike
anything in the training data.

\subsection{Workload Drift Detector Module}
\label{subsec:drift_detector}

The Workload Drift Detector Module, implemented in
\texttt{ml/drift\_detector.py}, implements the Page-Hinkley sequential
change-point test. The \texttt{PageHinkley} class maintains four state
variables: the observation count $n$, a running exponentially weighted mean
$\bar{x}$, the cumulative sum $S_n$, and the running minimum $m_n$. For each
new observation $x_t$, the update rule is:

\begin{equation}
  PH_n = \left[\sum_{t=1}^{n}(x_t - \bar{x}_t - \delta)\right] -
         \min_{1 \le t \le n}\left[\sum_{i=1}^{t}(x_i - \bar{x}_i - \delta)\right]
  \label{eq:pagehinkley}
\end{equation}

Drift is signalled when $PH_n > \lambda$. In QueryPilot, $\delta = 0.005$ and
$\lambda = 25$ (derived from the drift sensitivity parameter 0.5 as
$\lambda = \lfloor 50 \times (1 - 0.5) \rfloor$).

The \texttt{WorkloadDriftDetector} class wraps two \texttt{PageHinkley}
instances: one global detector across all queries, and one per-query detector
keyed by the first 60 characters of the query string with whitespace removed.
The \texttt{update(sql, execution\_time\_ms)} method updates both detectors
and returns the result of the global detector.

\subsection{SAGE Controller Module}
\label{subsec:sage_controller}

The SAGE Controller Module, implemented in \texttt{sage/controller.py}, is the
decision-making core of QueryPilot. The complete SAGE algorithm is:

\begin{enumerate}
  \item \textbf{Input:} SQL query string, MySQL cursor.
  \item Generate candidates using the Candidate Generator.
  \item Extract features using the Feature Extractor.
  \item Convert features to a fixed-order vector using \texttt{features\_to\_vector}.
  \item Evaluate \texttt{not\_ready}: True if the model is untrained or
        \texttt{queries\_seen $<$ n\_min}.
  \item \textbf{If not\_ready:} set decision = gather\_evidence, selected = native.
  \item \textbf{Elif drift\_active:} set decision = fallback, selected = native,
        clear training buffer, reset drift flag.
  \item \textbf{Else:} predict $(\mu_{native}, \sigma_{native})$ for native features.
        For each candidate $c$:
        \begin{itemize}
          \item Extract features for $c$ (reuse native features if hint\_type is native).
          \item Predict $(\mu_c, \sigma_c)$.
          \item Compute $u_c = \sigma_c / \mu_c$.
          \item Skip if $u_c > \tau_u = 0.8$.
          \item Skip if $\mu_c > \mu_{native} \times (1 + \tau_r = 0.15)$.
          \item Track the candidate with the lowest $\mu_c$.
        \end{itemize}
        If best candidate is non-native: decision = adapt.
        Else: decision = fallback.
  \item Execute the selected SQL using \texttt{\_run(sql, cursor)}.
  \item Update the drift detector with actual execution time.
  \item Increment \texttt{queries\_seen}; append (fvec, time) to training buffer.
  \item If buffer size $\geq$ n\_min: retrain all three ensemble models.
  \item Log the decision to the CSV file via \texttt{MetricsLogger}.
  \item \textbf{Return:} result dictionary with execution time, decision, reason,
        prediction mean, uncertainty, and drift flag.
\end{enumerate}

\subsection{Workload Generator Module}
\label{subsec:workload_generator}

The Workload Generator Module, implemented in \texttt{workload/generator.py}
and \texttt{workload/queries.py}, provides five TPC-H query templates (Q1, Q3,
Q5, Q6, Q10) and four scenario generators. Each template is parameterized with
a set of concrete parameter values drawn uniformly at random using a seeded
numpy random number generator, ensuring full reproducibility. The four
scenarios are: stable (uniform random mix of all five query types), sudden
shift (first 60\% scan-heavy, last 40\% join-heavy), gradual shift (probability
of join-heavy queries increases linearly from 0 to 1), and data skew (all Q3
queries with parameter distribution shifting from BUILDING to AUTOMOBILE at
70\% of the run).

\subsection{Experiment Runner Module}
\label{subsec:experiment_runner}

The Experiment Runner Module, implemented in \texttt{experiments/baseline.py},
executes the three-configuration interleaved comparison experiment. For each
query in the workload, all three configurations --- native MySQL, ML-Only, and
QueryPilot + SAGE --- execute the query back-to-back before moving to the next
query. This interleaved design controls for MySQL buffer pool warm-up bias,
ensuring that all three configurations observe the same buffer pool state for
each query. The module computes mean latency, median latency, P95 latency,
total execution time, SAGE regression rate (adapt decisions only), ML-Only
regression rate, and QueryPilot per-query overhead. Results are saved to JSON
files in the results directory.

\subsection{Dashboard Module}
\label{subsec:dashboard}

The Dashboard Module, implemented in \texttt{dashboard/app.py}, provides a
four-page Streamlit web application. The Overview page displays latency summary
tables and bar charts comparing the three configurations. The Query Timeline
page plots per-query execution times as line charts with red triangle markers
at drift detection events. The SAGE Decisions page shows a pie chart of
decision type frequencies and a table of the most recent SAGE decisions with
their reasons. The Experiment Comparison page shows grouped bar charts
comparing mean latency and regression rates across all four scenarios. The
dashboard reads from result JSON files in the experiments results directory and
allows the user to select a result file from the sidebar.
